% K=3; illustrate the generating draw J=2. J is hidden from the observer.
% The undirected blue connector marks a joint draw, not a causal relationship.
\begin{tikzpicture}[
  figure picture,x=1mm,y=1mm,
  every node/.style={font=\figurefont\fontsize{8}{10}\selectfont,text=FigureInk,inner sep=1.2pt},
  >={Latex[length=1.6mm,width=1.2mm]},
  rvnode/.style={draw=FigureStroke,line width=.7pt,rounded corners=1.8pt,minimum width=15mm,minimum height=8mm,fill=white},
  sampled/.style={rvnode,draw=FigureBlue,fill=FigureBlueFill}
]
  \path[use as bounding box] (0,0) rectangle (165,51);
  \node[sampled] (x) at (17,28) {$X$};
  \node[rvnode] (y1) at (69,28) {$Y_1$};
  \node[sampled] (y2) at (106,28) {$Y_2$};
  \node[rvnode] (y3) at (143,28) {$Y_3$};
  \draw[FigureBlue,line width=.8pt] (x.north) -- (17,40) -- (106,40) -- (y2.north);
  \node[anchor=south] at (61.5,41) {按 $P_{XY}$ 联合抽样};
  \node[text=FigureBlue,anchor=south] at (106,45) {$J=2$};
  \node at (17,17) {$X\sim P_X$};
  \node at (69,17) {$Y_1\sim P_Y$};
  \node[text=FigureBlue] at (106,17) {联合配对};
  \node at (143,17) {$Y_3\sim P_Y$};
  \node[font=\FigureNoteFont\fontsize{7.5}{9}\selectfont,text=FigureMuted] at (106,11) {其余候选彼此独立，也独立于配对};
  \draw[FigureGrid,line width=.5pt] (7,5)--(155,5);
  \node[anchor=north] at (81,3) {观察者只见 $W=(X,Y_1,Y_2,Y_3)$，不知道 $J$};
\end{tikzpicture}
